% !TEX program = lualatex
% Generated by scripts/build_module_practice.py - do not edit by hand.
\DocumentMetadata{
  pdfstandard=UA-2,
  pdfversion=2.0,
  lang=en-US,
  tagging=on
}
\documentclass[11pt]{article}
\usepackage[letterpaper, margin=0.8in, top=0.6in, bottom=0.6in]{geometry}
\usepackage{fontspec}
\setmainfont{Fira Sans}
\usepackage{amsmath}
\usepackage{amssymb}
% Fira Sans has no U+2032; an unmapped glyph fails PDF/UA-2 (.notdef)
\usepackage{newunicodechar}
\newunicodechar{′}{\ensuremath{'}}
\usepackage{booktabs}
\usepackage{longtable}
\usepackage{array}
\usepackage{calc}
\usepackage{etoolbox}
\usepackage{xcolor}
\usepackage{graphicx}
\definecolor{accent}{HTML}{BF5700}
% pandoc emits \tightlist on compact lists
\providecommand{\tightlist}{\setlength{\itemsep}{0pt}\setlength{\parskip}{0pt}}
\usepackage{titlesec}
\titleformat{\section}{\large\bfseries\color{accent}}{}{0pt}{}
\titlespacing*{\section}{0pt}{14pt}{4pt}
\titleformat{\subsection}{\normalsize\bfseries}{}{0pt}{}
\titlespacing*{\subsection}{0pt}{10pt}{2pt}
\usepackage{hyperref}
\hypersetup{pdftitle={ECON 201 Practice Problems: Firms as Price Setters}, pdfauthor={Div Bhagia}, hidelinks}
\pagestyle{plain}
\setlength{\parindent}{0pt}
\setlength{\parskip}{4pt}
\begin{document}
{\LARGE\bfseries\color{accent} Practice Problems: Firms as Price Setters}\hfill{\large ECON 201}
\vspace{4pt}


\section{Lecture 6: What does it cost to make a car? Scale and the cost of production}\label{lecture-6-what-does-it-cost-to-make-a-car-scale-and-the-cost-of-production}

\subsection{1. Denise's bakery}\label{denises-bakery}

Denise runs a small bakery. Rent and the ovens cost \$300 a day whether the bakery bakes anything or not, and each loaf costs \$2 in flour and labor.

\textbf{(a)} Write down the bakery's cost function, \(C(Q)\), where \(Q\) is loaves per day.

\textbf{(b)} Find the total cost, average cost, and marginal cost at 50, 100, and 300 loaves a day.

\textbf{(c)} At what output does average cost fall to \$3 a loaf?

\textbf{(d)} The landlord doubles the rent, so fixed costs rise to \$600 a day. Which of the three cost measures change, and how?

\subsection{2. Fixed or variable, and for how long}\label{fixed-or-variable-and-for-how-long}

\textbf{(a)} Denise's daily expenses are listed below. Which vary with the number of loaves baked, and which do not?

\begin{itemize}
\tightlist
\item
  Rent for the shop.
\item
  Flour, yeast, and butter.
\item
  Staff paid by the hour, with more hours scheduled on busy days.
\item
  The yearly fee for the bakery's website.
\item
  Electricity for the ovens.
\end{itemize}

\textbf{(b)} Denise's ovens can bake at most 300 loaves a day on normal shifts. To bake more, Denise has to pay staff overtime at a higher hourly wage. What happens to the marginal cost of a loaf beyond 300, and why does the book describe this as a short-run cost?

\subsection{3. Lena's screen-printing shop}\label{lenas-screen-printing-shop}

Lena prints T-shirts. Her cost function, in dollars per day, is \(C(Q) = 90 + Q^2/10\), where \(Q\) is shirts per day.

\textbf{(a)} What is the marginal cost per shirt when Lena goes from 10 to 15 shirts a day? From 15 to 20?

\textbf{(b)} Is Lena's marginal cost constant, like Beautiful Cars' was? What might make each extra shirt cost more than the last?

\subsection{4. Returns to scale}\label{returns-to-scale}

\textbf{(a)} Each firm below doubles all of its inputs. Classify its technology as increasing, constant, or decreasing returns to scale.

\begin{itemize}
\tightlist
\item
  A brewery's output rises from 100 barrels to 250.
\item
  A print shop's output rises from 100 posters to 200.
\item
  A consulting firm's output rises from 100 reports to 150.
\end{itemize}

\textbf{(b)} The print shop has constant returns to scale in production. Can its average cost still fall as it prints more posters? Explain.

\subsection{5. Multiple choice}\label{multiple-choice}

\medskip

\textbf{Q1.} A firm's cost function is \(C(Q) = 300 + 2Q\). Its marginal cost is:

a) \$300.

b) \$2.

c) \(300 / Q + 2\).

d) \$302.

\medskip

\textbf{Q2.} A cereal producer's cost function is \(C(Q) = 2Q\), where \(Q\) is pounds of cereal. Select all the statements that are correct.

a) There are no fixed costs of production.

b) The marginal cost of production is 2.

c) The producer's average cost falls with output.

d) For any quantity \(Q\), average cost and marginal cost are the same.

\medskip

\textbf{Q3.} Select all the statements that are correct.

a) With constant returns to scale, doubling all inputs doubles output.

b) With decreasing returns to scale, doubling all inputs more than doubles output.

c) With economies of scale, cost per unit falls as the firm expands its production.

d) With increasing returns to scale, tripling all inputs more than triples output.

\medskip

\textbf{Q4.} Which of the following is an example of network economies of scale?

a) A bakery buys flour in bulk at a discount.

b) A messaging app becomes more useful to each user as more of their friends join.

c) A carmaker spreads the cost of designing a model over more cars.

d) Doubling the material in a brewery's pipe more than doubles the liquid it can carry.

\medskip

\textbf{Q5.} A carmaker has fixed costs and a constant cost per car. Which statement about its average cost (AC) and marginal cost (MC) is correct?

a) AC is below MC at every output.

b) AC equals MC at every output.

c) AC is above MC at every output, and gets closer to MC as output rises.

d) AC rises with output while MC stays constant.

\section{Lecture 7: Who can get away with raising prices? Demand and elasticity}\label{lecture-7-who-can-get-away-with-raising-prices-demand-and-elasticity}

\subsection{1. Nadia's kayak rentals}\label{nadias-kayak-rentals}

Nadia rents kayaks by the day at a lake. Her demand curve is \(Q = 60 - 2P\), where \(Q\) is kayaks rented per day and \(P\) is the price of a rental in dollars.

\textbf{(a)} How many kayaks does Nadia rent at a price of \$20?

\textbf{(b)} Nadia raises her price from \$20 to \$21. How many kayaks does she rent now? Find the percentage change in price and in demand, and the price elasticity of demand.

\textbf{(c)} Find the elasticity for a \$1 price rise starting from \$25, and for a \$1 price rise starting from \$10. Use whichever formula you like. Where is demand elastic, and where is it inelastic?

\textbf{(d)} Nadia starts quoting her prices in cents, so her demand curve becomes \(Q = 60 - 0.02P\), with \(P\) in cents. Her price rises from 2,000 cents to 2,100 cents. Find the elasticity. Is it different from your answer to part (b)? Why or why not?

\subsection{2. Two coffee shops}\label{two-coffee-shops}

One coffee shop sits on a street with five other coffee shops. Another is the only place to buy coffee inside a large library. Which shop is likely to face the more elastic demand? Which one can raise its price without losing many customers?

\subsection{3. Flat and steep demand curves}\label{flat-and-steep-demand-curves}

True or false: a steeper demand curve always has a lower elasticity than a flatter one. Explain.

\subsection{4. Which demand is more elastic?}\label{which-demand-is-more-elastic}

Suppose the price of each good below rises by 10\%. In each pair, which good's demand is more elastic? Explain using one of the reasons from the lecture: close substitutes, a luxury or a necessity, how narrowly the market is defined, time to adjust, or share of spending.

\begin{itemize}
\tightlist
\item
  Movie tickets, or visits to the doctor.
\item
  A new laptop, or a tube of toothpaste.
\item
  A burger at a food court with ten other restaurants, or a ride on the only ferry to an island.
\item
  Coca-Cola, or all soft drinks.
\item
  Gasoline in the week after its price rises, or gasoline over the next five years.
\end{itemize}

\subsection{5. Beach parking}\label{beach-parking}

A county raises the price of its beach parking lot from \$20 to \$25 a day. The number of cars in the lot falls from 400 to 240 a day. Many drivers park on nearby streets instead, so the number of people visiting the beach falls by only 5\%.

\textbf{(a)} Find the elasticity of demand for spaces in the parking lot.

\textbf{(b)} Find the elasticity of demand for visits to the beach, using the same 25\% price rise.

\textbf{(c)} Why is demand for the lot so much more elastic than demand for visiting the beach?

\subsection{6. Multiple choice}\label{multiple-choice}

\medskip

\textbf{Q1.} The diagram shows two alternative demand curves, D and D′, for a product. Select all the statements that are correct.

\begin{center}
\includegraphics[width=5.00in,alt={Two parallel downward-sloping straight lines, price on the vertical axis from 0 to 12,000 dollars in steps of 2,000 and quantity on the horizontal axis from 0 to 100 in steps of 20. The lower line, D, runs from 6,000 dollars at zero quantity to zero price at 60 units. The upper line, D prime, runs from 10,000 dollars at zero quantity to zero price at 100 units.}]{img/core-q7-6.pdf}
\end{center}

a) On demand curve D, when the price is \$4,000, the firm can sell 30 units of the product.

b) On demand curve D′, the firm can sell 60 units at a price of \$4,000.

c) At a price of \$2,000, the firm can sell 40 more units of the product on D′ than on D.

d) With an output of 20 units, the firm can charge \$2,000 more on D′ than on D.

\medskip

\textbf{Q2.} A shop sells 20 hats per week at \$10 each. When it increases the price to \$12, the number of hats sold falls to 15 per week. Select all the statements that are correct.

a) When the price increases from \$10 to \$12, demand decreases by 25\%.

b) A 20\% increase in the price causes a 25\% fall in demand.

c) The demand for hats is inelastic.

d) Using these figures, we can estimate the elasticity of demand to be 1.25.

\medskip

\textbf{Q3.} The figure shows two straight-line demand curves, D₁ and D₂, which cross at point E. Select all the statements that are correct.

\begin{center}
\includegraphics[width=5.00in,alt={Two downward-sloping straight lines with no numbers on the axes, price on the vertical axis and quantity on the horizontal axis. D1 is steep and D2 is flatter, and they cross at point E. Point A is on D1 above and to the left of E. Point C is on D1 below and to the right of E. Point B is on D2 below and to the right of E, at a higher quantity and a higher price than C.}]{img/core-q7-8.pdf}
\end{center}

a) At point E, demand on D₁ is less elastic than demand on D₂.

b) Points A and C are on the same straight line, so demand is equally elastic at A and at C.

c) The firm facing D₁ probably has more close competitors than the firm facing D₂.

d) On D₂, demand is more elastic at E than at B.

\medskip

\textbf{Q4.} A firm raises its price by 4\%, and the quantity it sells falls by 8\%. The price elasticity of demand is:

a) 0.5, so demand is inelastic.

b) 2, so demand is elastic.

c) \(-2\), so demand is inelastic.

d) 2, so demand is inelastic.

\medskip

\textbf{Q5.} A campus concert raises its ticket price from \$20 to \$25, and ticket sales fall from 400 to 320. The price elasticity of demand is:

a) 0.8, so demand is inelastic.

b) 1.25, so demand is elastic.

c) 0.8, so demand is elastic.

d) 16, so demand is elastic.

\section{Lecture 8: Steak, ramen, and recessions: Other elasticities}\label{lecture-8-steak-ramen-and-recessions-other-elasticities}

\subsection{1. Devon's month}\label{devons-month}

Devon works at a warehouse. His hours go up, and his monthly income rises from \$2,000 to \$2,500. His purchases for the month change as below.

\begin{longtable}[]{@{}lrr@{}}
\toprule\noalign{}
& At \$2,000 & At \$2,500 \\
\midrule\noalign{}
\endhead
\bottomrule\noalign{}
\endlastfoot
Rides on the city bus & 40 & 30 \\
Gallons of milk & 8 & 9 \\
Rideshare trips & 2 & 3 \\
\end{longtable}

\textbf{(a)} Find the percentage change in demand and the income elasticity of demand for each good.

\textbf{(b)} Which of the three goods are normal, and which are inferior? Which of the normal goods is a necessity, and which is a luxury?

\textbf{(c)} The next year Devon's hours are cut and his income falls by 10\%. Using the elasticities from part (a), what happens to each of the three goods?

\subsection{2. Carla's grocery bill}\label{carlas-grocery-bill}

Ground beef at Carla's grocery store goes up from \$4.00 to \$5.00 a pound. Her purchases for the month change as below.

\begin{longtable}[]{@{}lrr@{}}
\toprule\noalign{}
& At \$4.00 & At \$5.00 \\
\midrule\noalign{}
\endhead
\bottomrule\noalign{}
\endlastfoot
Pounds of chicken & 10 & 15 \\
Packs of hamburger buns & 8 & 6 \\
Bottles of laundry detergent & 2 & 2 \\
\end{longtable}

\textbf{(a)} Find the percentage change in demand and the cross-price elasticity of demand for each good, with respect to the price of ground beef.

\textbf{(b)} Which of these goods are substitutes for ground beef, which are complements, and which is neither?

\subsection{3. Reading the elasticities}\label{reading-the-elasticities}

Suppose a survey of US households gives the five estimates below.

\begin{itemize}
\tightlist
\item
  The income elasticity of demand for visits to a laundromat is \(-0.4\).
\item
  The income elasticity of demand for bottled water is 0.3.
\item
  The income elasticity of demand for cruise vacations is 2.2.
\item
  The cross-price elasticity of demand for tea with respect to the price of coffee is 0.7.
\item
  The cross-price elasticity of demand for bagels with respect to the price of cream cheese is \(-0.8\).
\end{itemize}

\textbf{(a)} What does each number tell you about the good, or about the pair of goods?

\textbf{(b)} Incomes fall in a recession. Which of the three goods in the list sells better, and which loses the most demand?

\subsection{4. Omar's taco truck}\label{omars-taco-truck}

Omar runs a taco truck downtown and sells 500 tacos a week.

\textbf{(a)} The burrito place across the street raises its prices by 10\%, and Omar's sales rise from 500 to 520 tacos a week. Find the cross-price elasticity of demand for his tacos with respect to the price of a burrito, and say what it tells you.

\textbf{(b)} Later that year a warehouse nearby closes, and average income downtown falls by 8\%. Omar's sales rise from 500 to 520 tacos a week again. Find the income elasticity of demand for his tacos, and say what it tells you.

\subsection{5. Sorting goods}\label{sorting-goods}

\textbf{(a)} In each pair below, which good has the higher income elasticity of demand? Explain in a sentence.

\begin{itemize}
\tightlist
\item
  A monthly bus pass, or a plane ticket for a spring break trip.
\item
  Store-brand pasta, or dinner delivered from a restaurant.
\item
  Rent on a shared apartment, or a weekend at a hotel.
\end{itemize}

\textbf{(b)} For each pair below, say whether the cross-price elasticity of demand is positive, negative, or about zero, and whether the goods are substitutes, complements, or neither.

\begin{itemize}
\tightlist
\item
  Butter and margarine.
\item
  Printers and ink cartridges.
\item
  Gasoline and rides on the city bus.
\item
  Notebooks and lawn mowers.
\end{itemize}

\subsection{6. Multiple choice}\label{multiple-choice}

\medskip

\textbf{Q1.} A household's income rises by 10\%, and the number of cans of soup it buys falls by 4\%. The income elasticity of demand for soup is:

a) 0.4, so soup is a normal good.

b) \(-0.4\), so soup is an inferior good.

c) \(-2.5\), so soup is an inferior good.

d) \(-0.4\), so soup is a luxury.

\medskip

\textbf{Q2.} A family's monthly income rises from \$4,000 to \$5,000, and the number of restaurant meals it buys rises from 4 to 6 a month. The income elasticity of demand for restaurant meals is:

a) 0.5, so restaurant meals are a necessity.

b) 2, so restaurant meals are a luxury.

c) 2, so restaurant meals are a necessity.

d) 0.5, so restaurant meals are a luxury.

\medskip

\textbf{Q3.} Select all the statements that are correct.

a) An inferior good has a negative income elasticity of demand.

b) A good with an income elasticity of demand of 0.4 is a normal good and a necessity.

c) A luxury has an income elasticity of demand between 0 and 1.

d) Demand for an inferior good rises when income falls.

\medskip

\textbf{Q4.} The price of good B rises by 8\%, and demand for good A falls by 16\%. The cross-price elasticity of demand for A with respect to the price of B is:

a) 2, so A and B are substitutes.

b) \(-2\), so A and B are complements.

c) \(-0.5\), so A and B are complements.

d) \(-2\), so A and B are substitutes.

\medskip

\textbf{Q5.} For which pair of goods would you expect the cross-price elasticity of demand to be negative?

a) Coffee and tea.

b) Butter and margarine.

c) Hot dogs and hot dog buns.

d) Bus rides and rideshare trips.

\medskip

\textbf{Q6.} From the lecture: the income elasticity of demand in the United States is about 1.25 for recreation and about 0.35 for food, beverages, and tobacco. A recession cuts incomes by 10\%. Select all the statements that are correct.

a) Demand for recreation falls by about 12.5\%.

b) Demand for food, beverages, and tobacco falls by about 3.5\%.

c) Recreation is a necessity and food is a luxury.

d) Demand for an inferior good rises during the recession.

\section{Lecture 9: Sweet spot: Choosing the price and quantity that maximize profit}\label{lecture-9-sweet-spot-choosing-the-price-and-quantity-that-maximize-profit}

\subsection{1. Sierra Bikes}\label{sierra-bikes}

Sierra Bikes assembles electric bikes. Its demand curve is \(P = 1{,}200 - 10Q\), where \(Q\) is bikes per day and \(P\) is the price in dollars, and its cost function is \(C(Q) = 2{,}000 + 200Q\).

\textbf{(a)} Calculate Sierra's revenue, cost, and profit if it sells 20, 30, 40, 50, and 60 bikes a day.

\textbf{(b)} Which quantity gives the most profit, and what price does Sierra charge there?

\textbf{(c)} Now zoom in around 50 bikes. To sell one more bike, Sierra has to lower its price by \$10. Find its revenue at 49, 50, and 51 bikes. How much do the 50th and the 51st bikes each add to revenue? Marginal cost is \$200 a bike. Is each of those bikes worth selling?

\textbf{(d)} Sierra's landlord raises the rent, so its fixed cost rises from \$2,000 to \$3,000 a day. Marginal cost is still \$200 a bike. Does the profit-maximizing quantity or price change? What is Sierra's profit now?

\subsection{2. Marisol's cakes}\label{marisols-cakes}

Marisol bakes custom birthday cakes. From past orders she knows how many cakes she sells in a day at each price. Her cost function is \(C(Q) = 30 + 20Q\), a \$30 daily fee for a shared kitchen and \$20 of ingredients for each cake.

\begin{longtable}[]{@{}cc@{}}
\toprule\noalign{}
Price, \(P\) & Cakes sold, \(Q\) \\
\midrule\noalign{}
\endhead
\bottomrule\noalign{}
\endlastfoot
\$55 & 1 \\
\$50 & 2 \\
\$45 & 3 \\
\$40 & 4 \\
\$35 & 5 \\
\$30 & 6 \\
\end{longtable}

\textbf{(a)} Find her revenue, cost, and profit at each price. Which price gives the most profit?

\textbf{(b)} Each row of the table adds one cake, so the change in revenue from one row to the next is the marginal revenue of that cake. Find the marginal revenue of each cake. Her marginal cost is \$20 a cake. Which cakes are worth baking?

\subsection{3. Sorting}\label{sorting}

For each firm below, say whether it should sell more, sell less, or stay where it is.

\begin{enumerate}
\def\labelenumi{\roman{enumi}.}
\tightlist
\item
  A firm with \(\text{MR} = 40\) and \(\text{MC} = 25\) at its current quantity.
\item
  A firm with \(\text{MR} = 12\) and \(\text{MC} = 30\) at its current quantity.
\item
  A firm with \(\text{MR} = 18\) and \(\text{MC} = 18\) at its current quantity.
\end{enumerate}

\subsection{4. Multiple choice}\label{multiple-choice}

\medskip

\textbf{Q1.} A smoothie stand's demand curve is \(P = 10 - Q/20\), where \(Q\) is smoothies sold per day. Select all the statements that are correct.

a) To sell 100 smoothies a day, the highest price the stand can charge is \$5.

b) If the stand sets a price of \$6, it sells 80 smoothies a day.

c) The stand could sell 120 smoothies a day at a price of \$6.

d) Choosing how many smoothies to sell and choosing the price are the same decision.

\medskip

\textbf{Q2.} \emph{(Adapted from CORE Question 7.9)} Beautiful Cars has the cost function \(C(Q) = 60{,}000 + 10{,}000Q\). It maximizes its profit by producing \(Q^* = 20\) cars and setting the price at \(P^* = \$30{,}000\), for a profit of \$340,000. Suppose that the firm chooses instead to produce \(Q = 20\) cars and sets the price at \(P = \$29{,}800\). Which statement is correct?

a) The profit remains the same at \$340,000.

b) The profit is reduced to \$336,000.

c) The average cost of production is \$17,000.

d) The firm is unable to sell all the cars.

\medskip

\textbf{Q3.} \emph{(Adapted from CORE Question 7.10)} Beautiful Cars maximizes its profit at \(P^* = \$30{,}000\) and \(Q^* = 20\), and its marginal cost is \$10,000 for every car. Suppose that the firm decides to switch to a higher price. Which statement is correct?

a) The firm reduces the quantity of cars produced.

b) The marginal cost of producing an extra car is higher.

c) The total cost of production is higher.

d) The profit increases due to the new higher price.

\medskip

\textbf{Q4.} \emph{(Adapted from CORE Question 7.11)} The figure shows marginal cost, demand, and marginal revenue for Beautiful Cars. Select all the statements that are correct.

\begin{center}
\includegraphics[width=5.00in,alt={Three lines, price and cost on the vertical axis with marks at minus 10,000, 0, 10,000, 30,000, and 50,000 dollars, and cars per day on the horizontal axis with marks at 0, 10, 20, 25, 30, and 40. The demand curve falls from 50,000 dollars at zero cars. The marginal revenue line starts at the same point, falls twice as steeply, passes 30,000 dollars at 10 cars, crosses zero at 25 cars, and reaches minus 10,000 dollars at 30 cars. Marginal cost is a flat dashed line at 10,000 dollars. Marginal revenue crosses marginal cost at 20 cars, at the point labelled E prime, directly below the point E on the demand curve at 30,000 dollars.}]{img/core-q7-11.pdf}
\end{center}

a) At 10 cars, marginal revenue is above marginal cost, so selling one more car raises profit.

b) At 30 cars, marginal revenue is negative, so the firm's total profit must be negative.

c) The firm maximizes profit at 20 cars, where marginal revenue equals marginal cost, and charges \$30,000.

d) The firm should produce 25 cars, where marginal revenue is zero.

\medskip

\textbf{Q5.} A firm's marginal cost is \$4 at every quantity. The marginal revenue from one more unit is \$8 at 10 units, \$6 at 20, \$4 at 30, and \$2 at 40, falling steadily in between. Which quantity maximizes its profit?

a) 10

b) 20

c) 30

d) 40

\medskip

\textbf{Q6.} A firm sells 30 units at \$600 each. To sell a 31st unit it has to lower its price to \$590. What is its marginal revenue from the 31st unit?

a) \$10

b) \$290

c) \$590

d) \$600

\medskip

\textbf{Q7.} At its current quantity a firm has marginal revenue of \$200 and marginal cost of \$100. What should it do?

a) Sell more, because the next unit adds more to revenue than to cost.

b) Sell less, because marginal revenue is falling.

c) Stay where it is, because profit is already highest.

d) Raise the price and sell the same quantity.

\medskip

\textbf{Q8.} A firm's fixed cost rises by \$500 a day, and nothing else changes. Select all the statements that are correct.

a) Its profit-maximizing quantity falls.

b) Its marginal cost rises by \$500.

c) Its profit at every quantity falls by \$500.

d) It charges the same price as before.

\section{Lecture 10: The surplus from a sale and who captures it}\label{lecture-10-the-surplus-from-a-sale-and-who-captures-it}

\subsection{1. Sierra Bikes}\label{sierra-bikes}

Sierra Bikes assembles electric bikes. Its demand curve is \(P = 1{,}200 - 10Q\), where \(Q\) is bikes per day and \(P\) is the price in dollars, and its cost function is \(C(Q) = 2{,}000 + 200Q\). Sierra maximizes its profit by selling 50 bikes a day at \$700 each.

\textbf{(a)} How much is the 10th buyer willing to pay? What is the joint surplus from selling this buyer a bike, and how much of it goes to the buyer and how much to Sierra?

\textbf{(b)} Draw Sierra's demand curve and marginal cost. Mark the point where it sells 50 bikes at \$700, and shade consumer surplus and producer surplus. Find each, and Sierra's profit.

\subsection{2. Escape Hour}\label{escape-hour}

Escape Hour is the only escape room in town. The figure shows its demand curve and its marginal cost: each extra player costs \$20 in staff time and supplies. Escape Hour maximizes its profit at the point \(E\), selling 30 tickets a day at \$50 each.

\begin{center}
\includegraphics[width=5.00in,alt={Escape Hour's demand curve, a straight line falling from 80 dollars at zero tickets to zero at 80 tickets, with price and cost on the vertical axis marked every 10 dollars from 0 to 80, and tickets sold per day on the horizontal axis marked every 10 from 0 to 80. Marginal cost is a flat dashed line at 20 dollars. The point E is on the demand curve at 30 tickets and 50 dollars, with dashed guides to both axes.}]{img/escape-room.pdf}
\end{center}

\textbf{(a)} Find consumer surplus and producer surplus.

\textbf{(b)} Who captures more of the surplus, the players or Escape Hour? Why?

\textbf{(c)} Suppose Escape Hour lowered its price to \$40. How many tickets would it sell? What would happen to consumer surplus and producer surplus?

\subsection{3. Multiple choice}\label{multiple-choice}

\medskip

\textbf{Q1.} A demand curve is \(P = 90 - 2Q\). How much is the 15th buyer willing to pay?

a) \$15

b) \$60

c) \$75

d) \$90

\medskip

\textbf{Q2.} A buyer is willing to pay \$120 for a jacket. The store's marginal cost is \$50, and the price is \$80. How much surplus does the buyer get from the sale?

a) \$30

b) \$40

c) \$70

d) \$120

\medskip

\textbf{Q3.} A firm's demand curve is \(P = 60 - Q\) and its marginal cost is \$20 at every quantity. It sells 20 units at \$40, and its fixed cost is \$150. Select all the statements that are correct.

a) Consumer surplus is \$200.

b) Producer surplus is \$400.

c) Profit is \$250.

d) Consumer surplus equals the price times the quantity sold.

\medskip

\textbf{Q4.} A firm's producer surplus is \$5,000 a day, and its fixed cost is \$1,500 a day. What is its profit?

a) \$1,500

b) \$5,000

c) \$3,500

d) \$6,500

\medskip

\textbf{Q5.} A buyer is willing to pay \$1,100 for a bike that costs Sierra \$200 to make. Sierra cuts the price from \$700 to \$650. Select all the statements that are correct.

a) The joint surplus from the sale stays at \$900.

b) The buyer's gain rises from \$400 to \$450.

c) Sierra's gain rises by \$50.

d) The price cut moves \$50 of surplus from Sierra to the buyer.

\medskip

\textbf{Q6.} Which firm is most likely to capture a large share of the surplus from its sales?

a) The only ferry company serving an island.

b) One of many identical food trucks parked on the same street.

c) A gas station next to three other gas stations.

d) A seller on a website where hundreds of others sell the same phone case.

\section{Lecture 11: Deadweight loss and who captures the surplus}\label{lecture-11-deadweight-loss-and-who-captures-the-surplus}

\subsection{1. Sierra Bikes}\label{sierra-bikes}

Sierra Bikes assembles electric bikes. Its demand curve is \(P = 1{,}200 - 10Q\), where \(Q\) is bikes per day and \(P\) is the price in dollars, and its cost function is \(C(Q) = 2{,}000 + 200Q\). Sierra maximizes its profit by selling 50 bikes a day at \$700 each.

\textbf{(a)} How many buyers value a bike at more than its \$200 cost? Find the deadweight loss at Sierra's profit-maximizing point.

\textbf{(b)} How much is the 70th buyer willing to pay? Does this buyer buy a bike at \$700? Suppose Sierra sold this buyer one extra bike for \$350, and kept charging everyone else \$700. How much does each side gain? Is this a Pareto improvement?

\subsection{2. Escape Hour}\label{escape-hour}

Escape Hour is the only escape room in town. The figure shows its demand curve and its marginal cost: each extra player costs \$20 in staff time and supplies. Escape Hour maximizes its profit at the point \(E\), selling 30 tickets a day at \$50 each.

\begin{center}
\includegraphics[width=5.00in,alt={Escape Hour's demand curve, a straight line falling from 80 dollars at zero tickets to zero at 80 tickets, with price and cost on the vertical axis marked every 10 dollars from 0 to 80, and tickets sold per day on the horizontal axis marked every 10 from 0 to 80. Marginal cost is a flat dashed line at 20 dollars. The point E is on the demand curve at 30 tickets and 50 dollars, with dashed guides to both axes.}]{img/escape-room.pdf}
\end{center}

\textbf{(a)} How much is the 40th player willing to pay? Does this player buy a ticket at \$50? Suppose Escape Hour sold this player one ticket for \$30, and kept charging everyone else \$50. How much does each side gain? Is this a Pareto improvement?

\textbf{(b)} How many players value a ticket at more than its \$20 cost? Shade the deadweight loss and find its value.

\subsection{3. Multiple choice}\label{multiple-choice}

\medskip

\textbf{Q1.} Select all the statements that are correct.

a) If a Pareto improvement is possible, the outcome is not Pareto efficient.

b) A Pareto efficient outcome is always fair.

c) If one person gets all of a pizza and the others get none, the outcome can still be Pareto efficient.

d) Moving \$10 from one person to another is a Pareto improvement.

\medskip

\textbf{Q2.} A firm sets one price for all its buyers. Which change is a Pareto improvement?

a) The firm sells one extra unit, at a price between its marginal cost and the buyer's willingness to pay, to a buyer who was not buying, and keeps the price the same for everyone else.

b) The firm raises its price by \$1 on every unit.

c) The firm lowers its price by \$1 on every unit.

d) The firm stops selling to the buyer with the lowest willingness to pay.

\medskip

\textbf{Q3.} A firm that sets a single price produces less than the quantity where demand meets marginal cost. Why?

a) Its marginal cost rises as it produces more.

b) Its fixed cost is too high to produce more.

c) To sell one more unit, it has to lower the price on every unit it already sells.

d) Buyers with a low willingness to pay do not value the good at more than it costs.

\medskip

\textbf{Q4.} A firm's demand curve is \(P = 100 - Q\) and its marginal cost is \$20 at every quantity. It maximizes profit by selling 40 units at \$60. Select all the statements that are correct.

a) Consumer surplus is \$800.

b) Producer surplus is \$1,600.

c) The deadweight loss is \$1,600.

d) All gains from trade would be realized at 80 units.

\medskip

\textbf{Q5.} \emph{(Adapted from CORE Question 7.12)} Beautiful Cars maximizes its profit by selling 20 cars a day at \$30,000 each. Which statement is correct?

a) Consumer surplus is the sum, over all buyers, of the difference between each buyer's willingness to pay and the price.

b) Producer surplus equals the firm's profit.

c) Deadweight loss is the loss the firm incurs by not selling more cars.

d) All possible gains from trade are achieved when the firm chooses its profit-maximizing output and price.

\medskip

\textbf{Q6.} Which of these are examples of price discrimination? Select all the statements that are correct.

a) A museum charges students less than other visitors for the same ticket.

b) A pizza place charges more for a large pizza than a small one.

c) A store cuts the price of a jacket for every customer at the end of the season.

d) A car dealer bargains with each buyer, so two buyers pay different prices for the same car.

\section{Lecture 12: The markup rule and market power}\label{lecture-12-the-markup-rule-and-market-power}

\subsection{1. Canyon Kayaks}\label{canyon-kayaks}

Canyon Kayaks rents kayaks at a lake. Its demand curve is \(P = 80 - 2Q\), where \(Q\) is kayaks rented per day and \(P\) is the price in dollars. Each rental costs it \$20, so \(\text{MC} = 20\). Canyon Kayaks maximizes its profit by renting 15 kayaks a day at \$50 each.

\textbf{(a)} Find Canyon Kayaks' profit margin and its price markup.

\textbf{(b)} The slope of the demand curve is \(-2\). Find the price elasticity of demand at 15 kayaks and \$50.

\textbf{(c)} Check that Canyon Kayaks follows the markup rule.

\textbf{(d)} A second kayak rental opens at the lake. What happens to the elasticity of Canyon Kayaks' demand, and to its markup?

\subsection{2. Phone cases}\label{phone-cases}

A company sells phone cases. Each case costs it \$15 to make, so \(\text{MC} = 15\). At the price it chooses, the price elasticity of demand for its cases is 4.

\textbf{(a)} Use the markup rule to find its markup and its price.

\textbf{(b)} A rival brand starts selling very similar cases, and the elasticity rises to 6. Find the new price.

\subsection{3. Pixel Forge}\label{pixel-forge}

Pixel Forge makes a video game. Writing and testing the game costs \$60 million, whatever the number of players. Each extra download then costs \$2 in server fees, so \(\text{MC} = 2\). Games of this kind sell for \$5 a download.

\textbf{(a)} Find the average cost per download at 10 million, 20 million, and 60 million downloads. Why does it fall?

\textbf{(b)} How many downloads does Pixel Forge need to break even at \$5? What happens to its profit if it sells fewer, or more?

\textbf{(c)} At \$5, 60 million people buy a game of this kind. Suppose several studios make similar games and split those buyers equally. What is each studio's profit if there is one studio? Three? Six? How many studios can this market support?

\subsection{4. Multiple choice}\label{multiple-choice}

\medskip

\textbf{Q1.} \emph{(Adapted from CORE Question 7.13)} Select all the statements that are correct. A firm's market power would increase if:

a) Another firm that sold a close substitute goes out of business.

b) Buyers in its market start to care more about price than about the features of each product.

c) The firm gets a patent on its product.

d) Buyers become more loyal to the firm's brand.

\medskip

\textbf{Q2.} A firm maximizes its profit by charging \$50. At that price, the price elasticity of demand for its product is 5. What is its marginal cost?

a) \$10

b) \$40

c) \$45

d) \$50

\medskip

\textbf{Q3.} Which statement describes a natural monopoly?

a) A firm that holds a patent on its product.

b) A firm whose product has no close substitutes because of its brand.

c) A market where one firm can supply all buyers at a lower average cost than two firms could.

d) A market where firms agree with each other to charge the same price.

\medskip

\textbf{Q4.} A city has hundreds of hair salons. Each one offers a slightly different style and service, and a new salon can open easily. What is the market structure?

a) Perfect competition

b) Monopolistic competition

c) Oligopoly

d) Monopoly

\medskip

\textbf{Q5.} \emph{(Adapted from CORE Question 7.15)} A large retail chain plans to open a superstore in a small town. Select all the statements that are correct.

a) Local shops argue that many of the chain's goods are close substitutes for theirs, so the chain faces inelastic demand for those goods.

b) The chain argues that many of its goods are close substitutes for those in local shops, so demand is elastic and competition will keep prices low.

c) Local shops argue that once they are driven out, the chain will face less competition and can raise its prices.

d) The chain argues that local shops sell goods so different from its own that their demand is highly elastic.

\medskip

\textbf{Q6.} A coffee shop starts an advertising campaign that builds loyalty among its customers. According to the markup rule, what happens?

a) Its demand becomes more elastic, so it lowers its markup.

b) Its demand becomes less elastic, so it raises its markup.

c) Its marginal cost falls, so it lowers its price.

d) Its demand curve is unchanged, because advertising does not affect buyers.

\medskip

\textbf{Q7.} Four companies serve almost all US wireless customers, each sells a somewhat different plan, building a network costs billions, and each company watches what the others charge. What is the market structure?

a) Perfect competition

b) Monopolistic competition

c) Oligopoly

d) Monopoly

\medskip

\textbf{Q8.} A city's water utility is a natural monopoly. Which response is most likely to keep water affordable?

a) Split the utility into four companies, each with its own pipes.

b) Leave one utility in place and regulate the price it can charge.

c) Let the utility set whatever price it chooses.

d) Require every home to pay the marginal cost of its water and nothing more.

\medskip

\textbf{Q9.} Training an AI chatbot costs a company billions of dollars before anyone uses it. Answering one more question then costs a few cents in computing. Select all the statements that are correct.

a) The company's average cost per question falls as more people use the chatbot.

b) Its average cost is above its marginal cost at every number of questions.

c) If it charged each user only the computing cost of their questions, it would break even.

d) A new company with few users would have a higher average cost than an established one with many, so it would struggle to match the established company's price.

\medskip

\textbf{Q10.} Which of these are ways a firm can reduce the competition it faces? Select all the statements that are correct.

a) Buying one of its rivals.

b) Agreeing with its rivals to raise prices.

c) Keeping rivals' products off a platform it controls.

d) Cutting its price because its marginal cost fell.

\section{Lecture 13: Setting prices with few rivals}\label{lecture-13-setting-prices-with-few-rivals}

Answer each question before you open its \textbf{Solution}.

\subsection{Multiple choice}\label{multiple-choice}

\medskip

\textbf{Q1.} In which market is a firm's price most likely to change what its rivals charge?

a) A big city with hundreds of clothing shops.

b) A shopping mall with three fast food places.

c) Thousands of farmers selling wheat into a world market.

d) A town whose only water supplier is a single utility.

\medskip

\textbf{Q2.} In Lectures 9 to 12, a firm chose its price from its demand curve. What is different when it has only one or two rivals?

a) Its marginal cost now depends on its rivals' output.

b) It no longer faces a demand curve.

c) Its demand curve shifts when a rival changes price, so its best price depends on what it expects the rival to do.

d) It must charge the same price as its rivals by law.

\medskip

\textbf{Q3.} \emph{(Adapted from CORE Question 7.14)} Two rental shops on the same beach each choose a high or a low price. Select all the statements that are correct.

a) The more loyal their customers are, the more likely both shops are to set high prices.

b) When customers care mostly about price, competition pushes both shops' prices down.

c) Each shop's best price depends on what it expects the other shop to charge.

d) Each shop's best price is the same whatever the other shop charges.

\medskip

\textbf{Q4.} Two coffee shops on the same street start offering very different drinks and atmospheres, so each builds its own group of regulars. What happens?

a) Each shop's demand becomes less elastic, so both can charge higher prices.

b) Each shop's demand becomes more elastic, so both must charge lower prices.

c) Their marginal costs fall, so both lower their prices.

d) Nothing changes, because the shops still sell coffee.

\medskip

\textbf{Q5.} Two airlines fly the same route and compete their fares down. Who gains?

a) Both airlines.

b) Only the larger airline.

c) Their passengers.

d) Nobody.

\medskip

\textbf{Q6.} Two airlines flying the same route secretly agree to keep their fares high. Select all the statements that are correct.

a) This is a cartel.

b) Together, the two airlines act like a single monopoly.

c) Agreements like this are illegal in most countries.

d) Passengers benefit, because both airlines earn more.

\end{document}
